%%%%%%%%%%%%%%%%%%%%%%%%%%%%%%%%%%%%%%%%%%%%%%%%%%%%%%%%%%%%%%%%%%%%%%
% Lớp 10 | Véctơ (chưa xét tọa độ) | Tổng và hiệu của hai véctơ
% Dạng 0H5?2-3: Tìm điểm nhờ đẳng thức véctơ
% Số câu: 3  (sinh tự động bởi tach_cau.py; từ nay thầy sửa trực tiếp file này)
%%%%%%%%%%%%%%%%%%%%%%%%%%%%%%%%%%%%%%%%%%%%%%%%%%%%%%%%%%%%%%%%%%%%%%

%%%%%%%%%%%%%%%% ĐÚNG/SAI %%%%%%%%%%%%%%%%

% nguồn: MaHoaNH_DungSai_Lop10
\begin{ex}%[0H5N2-3]%[ma:b92d91d4]
	Cho tam giác $ABC$ và điểm $M$ sao cho $\vec{MA}-\vec{MB}-\vec{MC}=\vec{0}$. Hãy tìm các khẳng  định \textbf{đúng}.
	\choiceTF 
		{$M$ là trung  điểm của $BC$}
		{$M$ là trung điểm của $AB$} 
		{$M$ là trung điểm của  $AC$}
		{\True $ABMC$ là hình bình  hành}
	\loigiai{ 
		 $$\vec{MA}-\vec{MB}-\vec{MC}=\vec{0} \Leftrightarrow \vec{BA}=\vec{MC}.$$
		Suy ra $ABMC$  là hình bình hành.
	 }
\end{ex}

% nguồn: MaHoaNH_DungSai_Lop10
\begin{ex}%[0H5N2-3]%[ma:f77d7885]
	Cho hình bình hành $ABCD$. $M$ là điểm bất  kì sao cho $\overrightarrow{BM}=\overrightarrow{AB}+\overrightarrow{AD}$. Hãy tìm các khẳng định \textbf{đúng}.
	 \choiceTF
		{\True $M$ là đỉnh thứ tư của hình bình hành  $CABM$}
		 {\True $\vec{BM}=\vec{AC}$}
		{$M$ là đỉnh  thứ tư của hình bình hành $ABCM$}
		{$M$ là  đỉnh thứ tư của hình bình hành $DABM$}
	\loigiai{ 
		\begin{center}
			\begin{tikzpicture}[scale=1,font=\footnotesize,line join=round,line cap=round,>=stealth]
					\draw (0,0) coordinate (C) --++(60:2) coordinate  (B)--++(0:3) coordinate (A)--++(-120:2) coordinate (D)--(C);
					\coordinate (M) at ($(B)+(C)-(A)$);  \draw (B)--(M)--(C)--(A);
					 \foreach \x/\g in {A/90,B/90,C/-90,D/-90,M/-90}\fill[black] (\x) circle (1pt) +(\g:.3) node{$\x$};
			\end{tikzpicture}
		\end{center}
		\begin{itemize}
			\item  \[\overrightarrow{BM}=\overrightarrow{AB}+\overrightarrow{AD}\Leftrightarrow \overrightarrow{BM}=\overrightarrow{AC}.\]
			Vậy $M$ là đỉnh thứ tư của  hình bình hành $CABM$.
			\item $CABM$ là hình bình hành nên  $\vec{BM}=\vec{AC}$.
		\end{itemize}
	} 
\end{ex}

% nguồn: MaHoaNH_DungSai_Lop10
\begin{ex}%[0H5H2-3]%[ma:ffa0696e]
	Cho tam giác đều $ABC$ có cạnh bằng $a$.  Hãy tìm các khẳng định \textbf{đúng}.
	\choiceTF 
		{\True $|\overrightarrow{AB}+\overrightarrow{BC}|=a$} 
		 {$|\overrightarrow{AB}+\overrightarrow{AC}|=a\sqrt{2}$}
		{\True  $|\overrightarrow{AB}-\overrightarrow{AC}|=a$}
		{\True $|\overrightarrow{AB}+\overrightarrow{AC}|=a\sqrt{3}$} 
	 \loigiai{
		\begin{center}
			\begin{tikzpicture}[scale=1,font=\footnotesize,line join=round,line cap=round,>=stealth]
					\draw (0,0) coordinate (A) --++(60:3) coordinate (B)--++(-60:3) coordinate (C)--(A); 
					\coordinate (D) at ($(B)+(C)-(A)$); 
					\path (intersection of  A--D and C--B) coordinate (M);
					 \draw (A)--(B)--(D)--(C)
						(A)--(D); 
					\foreach \x/\g in {A/180,B/90,C/0,D/90,M/90}\fill[black] (\x)  circle (1pt) +(\g:.3) node{$\x$};
				\end{tikzpicture}
		\end{center}
		\begin{itemize}
			\item  $|\overrightarrow{AB}+\overrightarrow{BC}|=|\overrightarrow{AC}|=a$.
			 \item
			Ta có $\overrightarrow{AB}+\overrightarrow{AC}=\overrightarrow{AD}$  (với $D$ là điểm thỏa mãn $ABDC$ là hình bình hành).\\
			Gọi $M$ là giao điểm của $AD$ và $BC$ $\Rightarrow M$ là trung điểm của $BC$ và $AD$ (tính chất hình bình  hành).\\
			Xét  $\triangle ABC$, có
			\begin{eqnarray*}
				&&AB^2=AM^2+BM^2  \\
				 &\Rightarrow&AM^2=AB^2-BM^2=AB^2-\left(\dfrac{BC}{2}\right)^2a^2-\left(\dfrac{a}{2}\right)^2=\dfrac{3a^2}{4}\\
				 &\Rightarrow&AM=\dfrac{a\sqrt{3}}{2}.
			\end{eqnarray*}
			Vậy $AD=2AM=2 \cdot \dfrac{a  \sqrt{3}}{2}=a\sqrt{3}$.
			\item $|\overrightarrow{AB}-\overrightarrow{AC}|=|\overrightarrow{CB}|=a$. 
		\end{itemize}
	} 
\end{ex}
